\newcommand{\titulus}{\nomenFesti{S. Romualdi, Abbatis.}
\dies{Die 19. Iunii.}}
\newcommand{\oratio}{\pars{Oratio.}

\noindent Deus, qui per beátum Romuáldum in Ecclésia tua eremíticam vitam renovásti, concéde, ut, nosmetípsos abnegántes et Christum sequéntes, felíciter ad cæléstia regna mereámur ascéndere.

\pars{Pro pace in universo mundo.} \scriptura{Sir. 50, 25; 2 Esdr. 4, 20; \textbf{H416}}

\vspace{-4mm}

\antiphona{II D}{temporalia/ant-dapacemdomine.gtex}

\vfill

\noindent Deus, a quo sancta desidéria, recta consília et iusta sunt ópera: da servis tuis illam, quam mundus dare non potest, pacem; ut et corda nostra mandátis tuis dédita, et hóstium subláta formídine, témpora sint tua protectióne tranquílla.

\noindent Per Dóminum nostrum Iesum Christum, Fílium tuum, qui tecum vivit et regnat in unitáte Spíritus Sancti, Deus, per ómnia sǽcula sæculórum.

\noindent \Rbardot{} Amen.}
\newcommand{\invitatorium}{\pars{Invitatorium.}

\vspace{-4mm}

\antiphona{IV}{temporalia/inv-mirabileminsanctis.gtex}}
\newcommand{\hymnusmatutinum}{\pars{Hymnus}

\cuminitiali{IV}{temporalia/hym-LaetiColentes.gtex}}
\newcommand{\matversus}{\noindent \Vbardot{} Quam dúlcia fáucibus meis elóquia tua, Dómine.

\noindent \Rbardot{} Super mel ori meo.}
\newcommand{\lectioi}{\pars{Lectio I.} \scriptura{Iudc. 13, 1-14}

\noindent De libro Iúdicum.

\noindent In diébus illis: Rursum fílii Israel fecérunt malum in conspéctu Dómini, qui trádidit eos in manus Philisthinórum quadragínta annis.

\noindent Erat autem vir quidam de Sáraa et de stirpe Dan nómine Mánue habens uxórem stérilem.

\noindent Cui appáruit ángelus Dómini et dixit ad eam: «Ecce stérilis es et absque líberis, sed concípies et páries fílium. Cave ergo, ne vinum bibas ac síceram nec immúndum quidquam cómedas, quia ecce concípies et páries fílium, cuius non tanget caput novácula: erit enim puer nazarǽus Dei ex matris útero et ipse incípiet liberáre Israel de manu Philisthinórum».

\noindent Quæ cum venísset ad marítum, dixit ei: «Vir Dei venit ad me habens aspéctum sicut ángelus Dómini, terríbilis nimis. Non interrogávi eum, unde esset, nec ipse nomen suum mihi indicávit. Et dixit mihi: “Ecce concípies et páries fílium; cave, ne vinum bibas et síceram et ne áliquo vescáris immúndo: erit enim puer nazarǽus Dei ex útero matris usque ad diem mortis suæ”».

\noindent Orávit ítaque Mánue Dóminum et ait: «Obsecro, Dómine, ut vir Dei, quem misísti, véniat íterum et dóceat nos, quid debeámus fácere de púero, qui nascitúrus est».

\noindent Exaudivítque Deus precántem Mánue, et venit rursum ángelus Dei ad mulíerem sedéntem in agro. Mánue autem marítus eius non erat cum ea.

\noindent Festinávit ergo et cucúrrit ad virum suum nuntiavítque ei dicens: «Ecce appáruit mihi vir, qui illo die vénerat ad me».

\noindent Qui surréxit et secútus est uxórem suam veniénsque ad virum dixit ei: «Tu es, qui locútus es mulíeri?».

\noindent Et ille respóndit: «Ego sum».

\noindent Cui Mánue: «Quando, inquit, sermo tuus fúerit explétus, quid circa púerum observáre et fácere debémus?».

\noindent Dixítque ángelus Dómini ad Mánue: «Ab ómnibus, quæ locútus sum uxóri tuæ, abstíneat se; et quidquid ex vínea náscitur, non cómedat, vinum et síceram non bibat, nullo vescátur immúndo, et, quod ei præcépi, custódiat».}
\newcommand{\responsoriumi}{\pars{Responsorium 1.} \scriptura{\Rbardot{} Ex. 4, 13 \Vbardot{} Ps. 79, 2; \textbf{H417}}

\vspace{-5mm}

\responsorium{V}{temporalia/resp-obsecrodomine-CROCHU.gtex}{}}
\newcommand{\lectioii}{\pars{Lectio II.} \scriptura{Iudc. 13, 15-25}

\noindent Dixítque Mánue ad ángelum Dómini: «Obsecro, ut retineámus te et faciámus tibi hædum de capris».

\noindent Cui respóndit ángelus Dómini: «Si me rétines, non cómedam panes tuos; sin autem vis holocáustum fácere, offer illud Dómino».

\noindent Et nesciébat Mánue quod ángelus Dómini esset.

\noindent Dixítque ad eum: «Quod est tibi nomen, ut, si sermo tuus fúerit explétus, honorémus te?».

\noindent Cui ille respóndit: «Cur quæris nomen meum, quod est mirábile?».

\noindent Tulit ítaque Mánue hædum de capris et oblatiónem símilæ et pósuit super petram ófferens Dómino, qui facit mirabília; ipse autem et uxor eius intuebántur.

\noindent Cumque ascénderet flamma de altári in cælum, ángelus Dómini in flamma páriter ascéndit.

\noindent Quod cum vidísset Mánue et uxor eius, proni cecidérunt in terram; et ultra non eis appáruit ángelus Dómini.

\noindent Statímque intelléxit Mánue ángelum esse Dómini et dixit ad uxórem suam: «Morte moriémur, quia vídimus Deum».

\noindent Cui respóndit múlier: «Si Dóminus nos vellet occídere, de mánibus nostris holocáustum et oblatiónem non suscepísset nec ostendísset nobis hæc ómnia neque tália dixísset».

\noindent Péperit ítaque fílium et vocávit nomen eius Samson. Crevítque puer, et benedíxit ei Dóminus. Cœpítque spíritus Dómini impéllere eum in Castris Dan inter Sáraa et Esthaol.}
\newcommand{\responsoriumii}{\pars{Responsorium 2.} \scriptura{\Rbardot{} Lc. 1, 76 \Vbardot{} ibid., 77; \textbf{H82} \& \textbf{E269}}

\vspace{-5mm}

\responsorium{I}{temporalia/resp-tupuerpropheta-CROCHU.gtex}{}}
\newcommand{\lectioiii}{\pars{Lectio III.} \scriptura{Cap. 4 : PL 144,958-959}

\noindent E Vita sancti Romuáldi a sancto Petro Damiáno epíscopo conscrípta.

\noindent Cum in Romuáldi ánimo perfectiónis amor magis ac magis in dies crésceret, nullámque mens eius réquiem inveníret, audívit quia in Venetiárum pártibus quidam spiritális vir esset, Marínus nómine, qui eremíticam dúceret vitam.

\noindent Consénsu ab abbáte et frátribus nimírum facíllime impetráto, ad prædíctum venerábilem virum navígio discurrénte pervénit, et sub illíus regímine dégere humíllima mentis devotióne instítuit.

\noindent Erat autem Marínus, inter céteras virtútes, vir símplicis ánimi et sinceríssimæ ádmodum puritátis, nullo quidem magistério eremíticæ conversatiónis edóctus, sed ad hanc solúmmodo bonæ voluntátis impúlsu fúerat incitátus.

\noindent Pono autem hunc vivéndi morem habébat, ut per contínuum anni círculum tribus in hebdómada fériis médiam panis buccéllam et pugíllum fabæ gustáret, tribus vero vinum et pulméntum discréta sobrietáte percíperet. Psaltérium quidem per dies síngulos totum canébat.

\noindent Sed nimírum rudis, et nullátenus órdine vitæ singuláris instrúctus, sicut ipse póstmodum beátus Romuáldus hiláriter referébat, plerúmque, de cella éxiens, una cum discípulo passim per latitúdinem erémi psalléndo spatiabátur, nunc sub ista árbore vicénos, nunc sub illa tricénos, vel quadragénos cóncinens psalmos.

\noindent Romuáldus autem, quia sǽculum idióta relíquerat, apérto psaltério, vix suórum vérsuum notas syllabátim explicáre valébat; et hæc oculórum in ima defíxio intolerábilem sibi importunitátem acédiæ generábat; Marínus vero virgam in déxtera gerens, Romuáldo e divérso sedénti sinístram cápitis partem sæpíssime verberábat.

\noindent Post multa autem Romuáldus, gravi ádmodum necessitáte compúlsus, humíliter ait: «Magíster, a dextro me deínceps témpore pércute, quia iam lævæ auris audítum fúnditus perdo. »

\noindent Tunc ille, tantam eius patiéntiam admirátus, indiscrétæ severitátis témperat disciplínam.}
\newcommand{\responsoriumiii}{\pars{Responsorium 3.} \scriptura{\textbf{H373}}

\vspace{-5mm}

\responsorium{III}{temporalia/resp-istecognovit-CROCHU-cumdox.gtex}{}

\rubrica{vel ad libitum:} \pars{Responsorium 3.} \scriptura{\Rbardot{} Ps. 20, 3 \Vbardot{} ibid., 4; \textbf{H373}}

\vspace{-5mm}

\responsorium{IV}{temporalia/resp-desideriumanimaeeius-CROCHU-cumdox.gtex}{}}
\newcommand{\hymnuslaudes}{\pars{Hymnus}

\cuminitiali{VIII}{temporalia/hym-ORedemptorisPietas.gtex}}
\newcommand{\lectiobrevis}{\pars{Lectio Brevis.} \scriptura{Rom. 12, 1-2 }

\noindent Obsecro vos, fratres, per misericórdiam Dei, ut exhibeátis córpora vestra hóstiam vivéntem, sanctam, Deo placéntem, rationábile obséquium vestrum; et nolíte conformári huic sǽculo, sed transformámini renovatióne mentis, ut probétis quid sit volúntas Dei, quid bonum et bene placens et perféctum.}
\newcommand{\responsoriumbreve}{\pars{Responsorium breve.} \scriptura{Sir. 45, 9}

\antiphona{VI}{temporalia/resp-amaviteum.gtex}}
\newcommand{\benedictus}{\pars{Canticum Zachariæ.}

\vspace{-4mm}

\antiphona{IV e}{temporalia/ant-istecognovitiustitiam.gtex}

\vspace{-2mm}

\scriptura{Lc. 1, 68-79}

\vspace{-2mm}

\cantusSineNeumas
\initiumpsalmi{temporalia/benedictus-initium-iv-e2-auto.gtex}

%\vspace{-1.5mm}

\input{temporalia/benedictus-iv-e2.tex} \Abardot{}}
\newcommand{\preces}{\noindent Christum Deum sanctum, fratres, exaltémus,~\gredagger{} orántes ut serviámus illi in sanctitáte et iustítia coram ipso ómnibus diébus nostris,~\grestar{} et acclamémus:

\Rbardot{} Tu solus sanctus, Dómine.

\noindent Qui tentári voluísti per ómnia pro similitúdine nostra absque peccáto,~\grestar{} miserére nostri, Dómine Iesu.

\Rbardot{} Tu solus sanctus, Dómine.

\noindent Qui nos omnes ad perfectiónem caritátis vocásti,~\grestar{} sanctífica nos, Dómine Iesu.

\Rbardot{} Tu solus sanctus, Dómine.

\noindent Qui nos iussísti esse salem terræ et lucem mundi,~\grestar{} illúmina nos, Dómine Iesu.

\Rbardot{} Tu solus sanctus, Dómine.

\noindent Qui voluísti ministráre, non ministrári,~\grestar{} fac nos tibi et frátribus humíliter servíre, Dómine Iesu.

\Rbardot{} Tu solus sanctus, Dómine.

\noindent Tu, splendor glóriæ Patris et figúra substántiæ eius,~\grestar{} fac ut in glória vultum tuum respiciámus, Dómine Iesu.

\Rbardot{} Tu solus sanctus, Dómine.}
\newcommand{\benedicamuslaudes}{\cuminitiali{}{temporalia/benedicamus-memoria-laudes.gtex}}
\include{hebdomadaxi}
\include{feriavi}
